\section{Related Work}

\subsection{RL for LLM Reasoning}
Recent advanced~\citep{guo2025deepseek,zhu2024deepseek,openaio1,ouyang2022training} have demonstrated the effectiveness of reinforcement learning in LLM alignment and reasoning. Most existing methods train a reward model or use rule-based rewards for the entire response. In parallel, process supervision\citep{lightman2023letsverifystepstep} has shown a particularly promising performance than outcome supervision. 
Most existing works resort to training a process reward model (PRM) based on human- or auto-annotated process signals\citep{lightman2023letsverifystepstep, wang2024mathshepherdverifyreinforcellms,luo2024improvemathematicalreasoninglanguage,setlur2024rewardingprogressscalingautomated} and apply the static PRM for RL training~\citep{shao2024deepseekmathpushinglimitsmathematical,wang2024mathshepherdverifyreinforcellms}. However, a static PRM could suffer from distribution shift and reward hacking as RL training progresses. \citep{kazemnejad2024vineppounlockingrlpotential} overcomes this problem by conducting Monte Carlo rollouts to estimate the value for each step, but it suffers from high computation cost with around quadratic computational complexity and thus hinders scalability. 


\subsection{Tree Search for LLM Reasoning}
Tree search has mainly been explored in LLM alignment and inference, especially for data synthesis and offline preference training.
\citet{xie2024monte} employs MCTS to generate step-level preference pairs for DPO~\citep{rafailov2024direc}. ~\citet{chen2024alphamathzeroprocesssupervision,feng2024alphazeroliketreesearchguidelarge,zhang2024restmctsllmselftrainingprocess} employ MCTS to iteratively generate high-quality SFT data or produce process supervision signals for training. Other works explore reward-guided search~\citep{snell2024scalingllmtesttimecompute,yao2023treethoughtsdeliberateproblem,long2023largelanguagemodelguided} to boost the performance in the inference stage.  
However, few efforts have been devoted to studying how to explicitly integrate tree search into reinforcement learning training like AlphaZero~\citep{silver2017mastering} to improve LLM reasoning.












\hide{
\subsection{Tree Search for LLMs}
The integration of tree search methods into large language models has proven advantageous for both inference and training. 

\vpara{Inference Time Tree Search}Several studies have demonstrated the efficacy of tree search techniques at test-time. Notably, works such as \citep{long2023largelanguagemodelguided} and \cite{yao2023treethoughtsdeliberateproblem}, which implement Tree-of-Thought methods with depth-first and breadth-first searches, and \citep{hao2023reasoninglanguagemodelplanning}, which employs Monte Carlo Tree Search (MCTS), have shown that tree search at inference time, complemented by self-evaluation, enhances decision-making by exploring multiple reasoning paths. Furthermore, the combination of large language models with process reward-guided tree search has been shown to improve complex reasoning capabilities \cite{feng2024alphazeroliketreesearchguidelarge, park2024ensemblinglargelanguagemodels}. Additionally,\cite{snell2024scalingllmtesttimecompute} illustrates that reward model-guided inference-time tree search methods, such as beam search and lookahead search, serve as efficient test-time scaling techniques, often outperforming parameter scaling methods.

\vpara{Tree Search As Training Supervision Signal}Recent studies have demonstrated that by employing tree search mechanisms, researchers have successfully obtained high-quality supervisory signals that significantly enhance the performance of language models during the subsequent training phases. For instance, \citep{xie2024montecarlotreesearch} utilizes MCTS to generate step-level preference pairs, which serve as high-quality supervisory signals for preference learning in the Decision Process Optimization (DPO) stage. Additionally, several studies have employed tree searches to train policy-value models iteratively. In these frameworks, the policy model learns exclusively from high-quality sequences generated by MCTS, while the value model is trained on both positive and negative examples \citep{chen2024alphamathzeroprocesssupervision,feng2024alphazeroliketreesearchguidelarge,zhang2024restmctsllmselftrainingprocess}. These tree search methods have emerged as effective self-training techniques that substantially reduce annotation costs by autonomously generating high-quality supervisory signals. Collectively, these studies underscore the potential of tree search algorithms in improving the quality of training data for large language models (LLMs). Nonetheless, none have explicitly applied tree search within the context of reinforcement learning (RL) training procedures.

\subsection{Process Supervision}
Recent advancements in process supervision have yielded promising results, with process reward models (PRMs) demonstrating superior performance over traditional outcome reward models (ORMs) in various settings. This includes Best-of-N (BON) performance \citep{lightman2023letsverifystepstep, wang2024mathshepherdverifyreinforcellms}) and the ability of providing RL guidance \citep{wang2024mathshepherdverifyreinforcellms}. Despite the substantial utility of PRMs, their implementation incurs significant annotation costs, as each step's correctness must be meticulously labeled. To mitigate this issue, automated labeling methods have been developed. For example, \citep{luo2024improvemathematicalreasoninglanguage} and \citep{wang2024mathshepherdverifyreinforcellms} present approaches that estimate the value of current steps based on their probability of leading to future success. Other studies explore integrating process supervision signals into RL training. \citep{shao2024deepseekmathpushinglimitsmathematical} employs PRMs as the process supervision signal in RL, achieving marginal performance gains over ORMs. Additionally, \citep{setlur2024rewardingprogressscalingautomated} combines advantage estimation, a progress measurement, with value approximation to boost both test-time search and RL. To address the issue of inaccurate credit assignment by PRMs, \citep{kazemnejad2024vineppounlockingrlpotential} proposes using unbiased Monte Carlo-based estimates as process supervision signals instead of PRMs. However, this method faces efficiency and scalability challenges due to discarded rollout segments.
}







\hide{
REINFORCE is an essential approach for optimizing language model performance. The process involves sampling multiple sequences from the policy $\pi_{\theta}(y|x)$, where $x$ is the input context and $y$ is the generated sequence. Each sequence is then assigned a reward $R(y, x)$ to evaluate its quality, with a baseline $b$ subtracted to reduce variance in gradient estimates. The goal is to maximize the expected reward while incorporating a Kullback-Leibler (KL) divergence term to maintain proximity to a reference policy $\pi_{\text{ref}}(.|x)$. The combined optimization objective is given by:

\begin{equation}
\max_{\pi_{\theta}} \mathbb{E}_{x \sim \mathcal{D}, y \sim \pi_{\theta}(.|x)} \left[R_{\phi}(y, x) - \beta D_{\text{KL}}(\pi_{\theta}(.|x) \| \pi_{\text{ref}}(.|x))\right]
\end{equation}

The policy parameters $\theta$ are updated to maximize this objective, therefore enhancing positive outcomes and mitigating negative ones.
}


\hide{REINFORCE stands as a critical approach for optimizing model performance. 
The quintessential steps of the REINFORCE algorithm applied to language models involve the following pipeline:

\vpara{Sampling}
Initially, independent and identically distributed(i.i.d.) multiple samples are generated from the policy $\pi_{\theta}(y|x)$, where $y$ denotes the generated sequence and $x$ represents the input context. Sampling is critical in exploring the action space for diverse possible outputs.

\vpara{Reward Attribution}
Upon obtaining the generated samples, each sequence $y$ is associated with a reward $R(y, x)$ to measure the quality or correctness of the generated sequence.To reduce the variance of the gradient estimates without introducing bias, a baseline $b$ is subtracted from the reward. A common choice for this baseline is the average of all rewards corresponding to the same prompt.

\vpara{Optimization Objective}
The core goal is to maximize the expected reward while incorporating a Kullback-Leibler (KL) divergence term to ensure the trained policy $\pi_{\theta}(y|x)$ remains close to a reference policy $\pi_{\text{ref}}(.|x)$. The combined optimization objective thus becomes:

\begin{equation}
\max_{\pi_{\theta}} \underset {x \sim D, y \sim \pi_{\theta}(.|x)}{\mathbb{E}} \left[R_{\phi}(y, x) - \beta D_{\text{KL}}(\pi_{\theta}(.|x) \| \pi_{\text{ref}}(.|x))\right]
\end{equation}

Here, $r_{\phi}(x, y)$ represents the reward function, and $\beta$ is a weighting factor for the KL divergence term. 

\vpara{Policy Gradient Update}

After computing the reward and adjusting for the baseline, the policy parameters $\theta$ are updated by maximizing the optimization objective above. This ensures that the learning process amplifies positive outcomes while mitigating the impact of negative ones.
}